% Seminar kick-off: introduction and rules
% Build: `make slides` from the repository root (or `latexmk` in this directory).

\documentclass[aspectratio=169,11pt,t]{beamer}

\usepackage[T1]{fontenc}
\usepackage[utf8]{inputenc}
\usepackage[english]{babel}
\usepackage[sfdefault,scaled=0.95]{FiraSans}
\usepackage[scaled=0.85]{FiraMono}
\usepackage{microtype}
\usepackage{booktabs}
\usepackage{array}
\usepackage{tikz}
\usetikzlibrary{positioning,arrows.meta,calc}

\usetheme{seminar}

\hypersetup{colorlinks=true,urlcolor=semAccent,linkcolor=semInk}

\title[Diploma seminar kick-off]{From Data to Decisions:\\ Machine Learning in Finance and Business}
\subtitle{Bachelor's diploma seminar: introduction and rules}
\author{Michał Woźniak}
\institute{Faculty of Economic Sciences, University of Warsaw}
\date{Academic year 2026/27}

\newcommand{\repo}{\href{https://github.com/wne-uw-mjwozniak/wne-uw-seminars}{\texttt{github.com/wne-uw-mjwozniak/wne-uw-seminars}}}

\begin{document}

\begin{frame}[plain,noframenumbering]
  \titlepage
\end{frame}

\begin{frame}{Agenda}
  \tableofcontents
\end{frame}

% =============================================================================
\section{The seminar}
% =============================================================================

\begin{frame}{What this seminar is for}
  \framesubtitle{Course 2400-PL3SL339A \textbullet{} 3 ECTS per semester \textbullet{} pass/fail}
  You write a bachelor's thesis that uses \key{modern quantitative methods} (machine learning,
  statistical learning, econometrics) to answer a well-posed question in \key{finance or business}.

  \bigskip
  What I care about most:
  \begin{itemize}
    \item a research question that matters to someone, stated precisely,
    \item rigorous validation: honest baselines, out-of-sample testing, robust conclusions,
    \item results that anyone can reproduce from your repository,
    \item a coherent story: \emph{question $\rightarrow$ data $\rightarrow$ model $\rightarrow$ decision or conclusion}.
  \end{itemize}

  \bigskip
  \muted{Applying a method is the easy part. Interpreting and criticising the result is the thesis.}
\end{frame}

\begin{frame}{Thematic scope}
  \framesubtitle{Deliberately broad: any topic with a sound question and a rigorous quantitative method}
  \begin{columns}[T,onlytextwidth]
    \column{0.48\textwidth}
    \begin{itemize}
      \item Comparing quantitative methods on a financial or business problem
      \item Case studies: finance, insurance, e-commerce, logistics
      \item Probabilistic modelling and forecasting
    \end{itemize}
    \column{0.48\textwidth}
    \begin{itemize}
      \item Risk measurement and uncertainty
      \item Causal inference for business decisions and policy evaluation
      \item ML beyond classic analytics: behavioural economics, public policy
    \end{itemize}
  \end{columns}

  \bigskip
  \begin{block}{Where I can help most}
    The intersection of econometrics and machine learning: quantitative finance, risk modelling,
    time-series forecasting, probabilistic methods, customer analytics. I combine academic research
    with commercial projects in insurance, e-commerce and logistics.
  \end{block}
\end{frame}

% =============================================================================
\section{How we work}
% =============================================================================

\begin{frame}{Format: remote, one joint start, then your own track}
  \begin{center}
  \begin{tikzpicture}[
      node distance=9mm,
      stage/.style={draw=semRule,fill=semTint,rounded corners=2pt,minimum height=15mm,
                   text width=36mm,align=center,font=\small},
      arr/.style={-{Stealth[length=2.2mm]},semAccent,line width=0.9pt}]
    \node[stage] (a) {\key{Joint kick-off}\\ rules, tooling, topics};
    \node[stage,right=of a] (b) {\key{Topic and team}\\ concept note approved};
    \node[stage,right=of b] (c) {\key{Individual track}\\ weekly status, kanban};
    \draw[arr] (a) -- (b);
    \draw[arr] (b) -- (c);
  \end{tikzpicture}
  \end{center}

  \medskip
  \begin{itemize}
    \item All meetings are \key{remote, on Google Meet}. This is the preferred way to reach me.
    \item After the kick-off we work individually or in project teams, not as a lecture group.
    \item Joint sessions return at two points: \emph{research design} and \emph{interpreting results},
          where you present and get feedback from the group.
    \item Asynchronous questions go to \key{GitHub issues} in your project repository, so the
          answer stays next to the work. University e-mail is for formal matters.
  \end{itemize}
\end{frame}

\begin{frame}{Weekly status and a kanban board}
  \framesubtitle{Every project lives in a GitHub Project. The board is the single source of truth for progress.}
  \begin{center}
  \begin{tikzpicture}[
      col/.style={draw=semRule,fill=semTint,rounded corners=2pt,minimum width=25mm,minimum height=27mm},
      head/.style={font=\scriptsize\bfseries,text=semAccent},
      card/.style={draw=semRule,fill=white,rounded corners=1pt,minimum width=21mm,minimum height=5mm,
                   font=\tiny,text=semInk,inner sep=1pt}]
    \foreach \i/\name in {0/Backlog,1/Ready,2/In progress,3/In review,4/Done} {
      \node[col] (c\i) at (\i*2.75,0) {};
      \node[head,anchor=north] at ($(c\i.north)+(0,-0.8mm)$) {\name};
    }
    \node[card] at ($(c0.north)+(0,-8mm)$)  {robustness checks};
    \node[card] at ($(c0.north)+(0,-14.5mm)$) {write data chapter};
    \node[card] at ($(c0.north)+(0,-21mm)$) {ES backtests};
    \node[card] at ($(c1.north)+(0,-8mm)$)  {HAR-RV baseline};
    \node[card] at ($(c1.north)+(0,-14.5mm)$) {rolling-window CV};
    \node[card,draw=semAccent] at ($(c2.north)+(0,-8mm)$) {GARCH baseline};
    \node[card] at ($(c3.north)+(0,-8mm)$)  {data pipeline PR};
    \node[card] at ($(c4.north)+(0,-8mm)$)  {concept note};
    \node[card] at ($(c4.north)+(0,-14.5mm)$) {repo + \texttt{uv} env};
  \end{tikzpicture}
  \end{center}

  \begin{itemize}
    \item \key{Every week: a status check-in.} What moved, what is blocked, what is next.
    \item I read the board before we talk, so the call is about decisions, not reporting.
  \end{itemize}
\end{frame}

\begin{frame}{Kanban: the working rules}
  \begin{columns}[T,onlytextwidth]
    \column{0.5\textwidth}
    \textbf{Cards}
    \begin{itemize}
      \item One card = one GitHub issue = one piece of work you can finish in days, not weeks.
      \item Each card states its \key{definition of done}: a figure, a table, a section, a passing test.
      \item Writing is work too: chapters and sections are cards.
    \end{itemize}
    \column{0.46\textwidth}
    \textbf{Flow}
    \begin{itemize}
      \item \key{Limit work in progress}: at most two cards \emph{In progress} per person.
      \item Work on branches. Open a pull request linked to the issue, then manually move its card to \emph{In review}.
      \item I review there. Merge into the default branch to close the issue and mark it \emph{Done}.
      \item \must{No movement on the board for weeks is the main risk signal} for passing the semester.
    \end{itemize}
  \end{columns}
\end{frame}

\begin{frame}{The toolchain is part of the course}
  \framesubtitle{These are learning outcomes in the syllabus, not preferences}
  \begin{description}[\quad Environment]
    \item[Writing] \must{LaTeX only} (Overleaf or a local TeX distribution); bibliography in BibTeX.
    \item[Versioning] Everything in \key{git}, hosted on GitHub: thesis sources and analysis code,
          with a readable commit history.
    \item[Language] \key{Python} preferred; R accepted.
    \item[Environment] Python: \texttt{uv} with \texttt{pyproject.toml}. R: \texttt{renv}.
    \item[Standard] \must{Reproducibility}: anyone with access to the repository can recreate your
          results without contacting you.
  \end{description}

  \bigskip
  \muted{New to some of these? That is expected. The first weeks include time to set them up, and
  the seminar repository has a thesis template and links to get started.}
\end{frame}

\begin{frame}[fragile]{What your project repository should look like}
  \begin{columns}[T,onlytextwidth]
    \column{0.44\textwidth}
\begin{semiverbatim}\footnotesize
my-thesis/
  README.md       \muted{how to reproduce}
  pyproject.toml  \muted{uv environment}
  uv.lock
  data/           \muted{raw data, fetch script}
  src/            \muted{reusable code}
  experiments/    \muted{configs, run scripts}
  results/        \muted{generated output}
  paper/          \muted{LaTeX sources, .bib}
\end{semiverbatim}
    \column{0.53\textwidth}
    \begin{itemize}
      \item One command rebuilds the results; one command rebuilds the PDF.
      \item Figures and tables are \key{generated by code}, never pasted by hand.
      \item Fix random seeds; record data versions and download dates.
      \item Data you cannot share: document how to obtain it and commit a small synthetic sample.
      \item No secrets, API keys or personal data in git, ever.
    \end{itemize}
  \end{columns}
\end{frame}

% =============================================================================
\section{Timeline and passing}
% =============================================================================

\begin{frame}{Indicative timeline}
  \framesubtitle{Exact submission and defence dates follow the WNE academic calendar}
  \begin{center}
  \begin{tikzpicture}[x=1.44cm,
      tick/.style={semMuted,line width=0.6pt},
      lab/.style={font=\scriptsize,text=semMuted},
      up/.style={font=\scriptsize,align=center,text width=23mm,anchor=south},
      down/.style={font=\scriptsize,align=center,text width=23mm,anchor=north}]
    \draw[semAccent,line width=1.4pt] (0,0) -- (3.55,0);
    \draw[semAccent!45,line width=1.4pt] (3.55,0) -- (8.3,0);
    \foreach \x/\m in {0/Oct,1/Nov,2/Dec,3/Jan,4/Feb,5/Mar,6/Apr,7/May,8/Jun} {
      \draw[tick] (\x,0.1) -- (\x,-0.1);
      \node[lab,anchor=north] at (\x,-0.12) {\m};
    }
    \node[up]   at (0,0.25)   {Kick-off, topic, team};
    \node[down] at (1,-0.62)  {Concept note approved};
    \node[up]   at (2,0.25)   {Data, repo, first baseline};
    \node[down] at (3,-0.62)  {\must{Semester 1 pass}: outline, questions, literature started};
    \node[up]   at (4.5,0.25) {Modelling and experiments};
    \node[down] at (6,-0.62)  {Full draft; \must{topic submitted for approval}};
    \node[up]   at (7,0.25)   {Draft read, revisions};
    \node[down] at (8,-0.62)  {Submission and defence};
  \end{tikzpicture}
  \end{center}

  \medskip
  \begin{itemize}
    \item Give me \key{time to read}: a full draft needs weeks, not days, before any deadline.
    \item The schedule is tight only if modelling starts in spring. Start in December.
  \end{itemize}
\end{frame}

\begin{frame}{What it takes to pass each semester}
  \begin{columns}[T,onlytextwidth]
    \column{0.48\textwidth}
    \begin{block}{After semester 1}
      \begin{itemize}
        \item thesis outline in LaTeX,
        \item data collected,
        \item research questions and hypotheses,
        \item literature review started,
        \item git repository with thesis sources and analysis code.
      \end{itemize}
    \end{block}
    \column{0.48\textwidth}
    \begin{block}{After semester 2}
      \begin{itemize}
        \item literature review and modelling complete,
        \item thesis ready for defence,
        \item code reproduces all results from the repository,
        \item for high-potential work: a draft of a scientific article.
      \end{itemize}
    \end{block}
  \end{columns}

  \bigskip
  In both semesters: \must{documented progress} (commits, successive versions of the document)
  and \must{active participation} in consultations and joint discussions.
  \src{course syllabus, USOS 2400-PL3SL339A}
\end{frame}

\begin{frame}{Your topic must be formally approved}
  \framesubtitle{Procedure of the WNE Teaching Council (Rada Dydaktyczna)}
  \begin{enumerate}
    \item We agree the topic together. I am responsible for checking that it fits your programme,
          level of study and leading discipline.
    \item I submit it to the Dean's Office during the last semester of the seminar,
          \must{no later than 4 weeks before you plan to submit the thesis}.
    \item The Teaching Council reviews the topic within 14 days. It may request changes,
          after which the corrected topic is submitted again.
    \item \must{A thesis without an approved topic cannot be entered into APD or defended.}
  \end{enumerate}

  \bigskip
  \muted{In practice: we fix the final title early in spring, and we word it so the economic or
  financial problem is visible in the title itself.}
  \src{Komunikat nr 1/2025 of the Chair of the WNE Teaching Council, 20 May 2025}
\end{frame}

% =============================================================================
\section{Faculty rules in brief}
% =============================================================================

\begin{frame}{What WNE expects from a bachelor's thesis}
  \begin{itemize}
    \item A \key{clearly defined research goal} and a research question or hypothesis.
    \item Evidence that you can do research using current domestic or international scholarly literature.
    \item Research tools consistent with your programme, or going beyond it.
    \item Either your own study, or an in-depth literature-based report. In this seminar we do
          \key{empirical work}.
    \item Traditional form: up to about \key{50 pages} (100,000 characters with spaces).
          Concise beats long.
    \item On a Polish-language programme the thesis may be written in \key{Polish or English}.
    \item It may take the \key{traditional form or the form of a scientific article}.
  \end{itemize}
  \src{\href{https://www.wne.uw.edu.pl/download_file/6305/254}{Resolution no.\ 7, consolidated text of 3 November 2025}}
\end{frame}

\begin{frame}{The method must serve an economic or financial problem}
  \framesubtitle{How the Teaching Council judges whether a topic fits the programme}
  \footnotesize
  \begin{tabular}{@{}>{\raggedright}p{0.24\textwidth}>{\raggedright\arraybackslash}p{0.72\textwidth}@{}}
    \toprule
    \textbf{Programmes} & \textbf{The thesis should} \\
    \midrule
    Economics \muted{(Committee 1)} &
      analyse economic phenomena at micro or macro scale, build on economic theory or policy, and
      pose a research problem grounded in the literature. \\
    \addlinespace
    Finance, accounting, quantitative finance \muted{(Committee 2)} &
      address a genuine finance question (pricing, risk management, portfolio optimisation, \ldots).
      Statistical, econometric and ML methods are welcome \emph{as long as they serve the financial
      problem}; a generic ``business'' theme is not enough. \\
    \addlinespace
    Informatics and Econometrics, data science \muted{(Committee 3)} &
      show independent, critical use of advanced quantitative methods, state the economic, social
      or financial context, and interpret results in those terms. \\
    \bottomrule
  \end{tabular}

  \medskip
  \normalsize
  \key{For us:} ``I benchmarked five models'' is not a topic. ``Which models forecast portfolio
  tail risk reliably enough for capital decisions'' is.
  \src{Komunikat nr 1/2026 of the Chair of the WNE Teaching Council, 23 March 2026}
\end{frame}

\begin{frame}{Formal requirements and the thesis template}
  \begin{columns}[T,onlytextwidth]
    \column{0.5\textwidth}
    \textbf{Required order}
    \begin{enumerate}\small
      \item Title page
      \item Signed declarations: supervisor and author(s)
      \item Abstract (max.\ 800 characters), keywords, Erasmus code, title in the other language
      \item Table of contents (max.\ three levels)
      \item Main text
      \item Bibliography
      \item Lists: abbreviations, tables, figures, appendices
    \end{enumerate}
    \column{0.46\textwidth}
    \textbf{Typesetting}
    \begin{itemize}\small
      \item Times-like font, 12 pt; line spacing 1.5
      \item All margins 25 mm; page numbers bottom centre
      \item Tables and figures numbered through the whole thesis, each with a title and a \key{source}
      \item Citations as (Author, year); bibliography alphabetical, preferably \key{APA}
    \end{itemize}
  \end{columns}
  \src{Załącznik B and Załącznik D to Resolution no.\ 7 of the WNE Teaching Council}
\end{frame}

\begin{frame}{From finished thesis to defence}
  \begin{enumerate}
    \item \key{Submit} the thesis file by e-mail to the Dean's Office; they open your entry in APD.
    \item \key{Upload} a searchable PDF to APD with the English title, abstract and keywords, then
          send it for the supervisor's approval.
    \item \key{Clear your accounts}: IT lab, WNE library, BUW, language and PE units.
    \item \key{Anti-plagiarism check} in JSA, run by the supervisor; I accept the thesis after
          reviewing the report.
    \item A \key{reviewer} is appointed. Both reviews are uploaded to APD at least 3 days before
          the defence.
    \item One \key{printed copy}, double-sided, soft-bound, with the AI declaration, delivered no
          later than 2 days before the exam.
  \end{enumerate}
  \medskip
  \muted{For June and September defences the submission deadline is strict to the hour (12:00).}
  \src{wne.uw.edu.pl \textrightarrow{} Prace dyplomowe \textrightarrow{} Składanie prac; Załącznik C}
\end{frame}

\begin{frame}{How the thesis is graded}
  \begin{columns}[T,onlytextwidth]
    \column{0.56\textwidth}
    \small
    \begin{tabular}{@{}lp{0.78\textwidth}@{}}
      \toprule
      \textbf{5!} & outstanding: well beyond the programme, a real contribution to science \\
      \textbf{5}  & very good: makes significant use of the programme's content, fully meets requirements \\
      \textbf{4}  & good: moderate use of the programme's content \\
      \textbf{3}  & weak: sufficient use, requirements partly met \\
      \textbf{2}  & fail: substantive errors leading to wrong conclusions \\
      \bottomrule
    \end{tabular}
    \column{0.4\textwidth}
    \begin{itemize}
      \item Supervisor and reviewer each score the thesis on the official review form;
            half grades are allowed.
      \item \key{Final grade = arithmetic mean} of the two.
      \item The reviewer sees only the thesis. Everything that matters has to be on the page.
    \end{itemize}
  \end{columns}
  \src{wne.uw.edu.pl \textrightarrow{} Prace dyplomowe \textrightarrow{} Ocena prac dyplomowych}
\end{frame}

% =============================================================================
\section{Think bigger: articles and teams}
% =============================================================================

\begin{frame}{The thesis can be a scientific article}
  \framesubtitle{This is the route I encourage for ambitious projects}
  \begin{itemize}
    \item The article has the structure typical for its field, about \key{10--12 thousand words},
          formatted for a chosen journal.
    \item A \key{supplement} is mandatory: literature review, methodology and results in more depth.
          It guides the reviewer and documents your own contribution.
    \item Article and supplement may be in different languages.
    \item Publication before the defence is \key{not required}.
    \item The supervisor may co-author the article; the supplement is the student's own work.
  \end{itemize}

  \smallskip
  \begin{block}{Layout of an article-form thesis}
    \small Title page \textrightarrow{} declarations \textrightarrow{} abstract page \textrightarrow{} contents
    \textrightarrow{} ``Scientific article'' \textrightarrow{} article \textrightarrow{} ``Supplement''
    \textrightarrow{} supplement \textrightarrow{} lists
  \end{block}
  \src{wne.uw.edu.pl \textrightarrow{} Prace w formie artykułów naukowych; Załącznik B, point 17}
\end{frame}

\begin{frame}{Co-authorship: what the rules allow}
  \begin{center}
  \begin{tabular}{@{}lccc@{}}
    \toprule
    \textbf{Form} & \textbf{Students} & \textbf{Each student} & \textbf{Supervisor} \\
    \midrule
    Traditional thesis   & 2 & exactly 50\%     & -- \\
    Scientific article   & 1 & at least 60\%    & at most 40\% \\
    Scientific article   & 2 & at least 40\%    & at most 20\% \\
    \bottomrule
  \end{tabular}
  \end{center}

  \medskip
  \begin{itemize}
    \item A thesis can have \must{at most two student authors}.
    \item Every author declares their substantive and percentage contribution, chapter by chapter.
    \item I must justify the co-authorship in my review.
  \end{itemize}
  \src{WNE thesis guidelines; declarations in Załącznik D}
\end{frame}

\begin{frame}{Larger projects, several theses, one paper}
  \begin{columns}[T,onlytextwidth]
    \column{0.5\textwidth}
    \textbf{The model}
    \begin{itemize}
      \item A research programme bigger than one thesis, split into \key{clearly separable parts}.
      \item Each part is one thesis, by one student or a pair, with its own question and contribution.
      \item Shared infrastructure: one repository or organisation, common data pipeline,
            common evaluation protocol.
      \item The parts combine into a \key{journal submission}.
    \end{itemize}
    \column{0.46\textwidth}
    \textbf{Why bother}
    \begin{itemize}
      \item A stronger benchmark than any one person can build.
      \item You learn to work the way research and industry teams do: issues, reviews, shared code.
      \item A paper and a public repository are worth more on a CV than a grade.
    \end{itemize}
    \medskip
    \muted{\small Authorship follows contribution, agreed and written down at the start of the project.}
  \end{columns}
\end{frame}

% =============================================================================
\section{AI tools and academic integrity}
% =============================================================================

\begin{frame}{AI tools: allowed, declared, and your responsibility}
  \begin{block}{Faculty requirement}
    Every thesis includes a signed \key{declaration on the use of AI systems}: which tool you used
    and why, or that you used none. You confirm that you verified the content and take full
    responsibility for the final thesis.
  \end{block}

  \medskip
  \textbf{Working rules in this seminar}
  \begin{itemize}
    \item Use AI like a capable colleague: code, debugging, language, critique. Not to produce
          the research idea, the interpretation or the conclusions.
    \item \must{You must be able to explain every line of code and every sentence.} I will ask.
    \item \must{Verify every reference.} Citing a paper that does not exist is a disqualifying error.
    \item Keep a short log of what you used AI for; the declaration then writes itself.
  \end{itemize}
\end{frame}

% =============================================================================
\section{Topics}
% =============================================================================

\begin{frame}{How to propose a topic}
  \begin{columns}[T,onlytextwidth]
    \column{0.46\textwidth}
    \textbf{First message: three things}
    \begin{enumerate}
      \item A working title.
      \item Two or three sentences: whose problem is it, what do they want to achieve, what is
            missing today?
      \item Confirmation that \key{the data exists} and you can get it.
    \end{enumerate}
    \medskip
    Your own idea is welcome. So is joining one of the projects on the next slides.
    \column{0.5\textwidth}
    \textbf{Then: a two-page concept note}
    \begin{enumerate}
      \item The problem
      \item Why it matters to an economist or a financial decision-maker
      \item Research questions and hypotheses
      \item Methods and data
      \item Planned structure
    \end{enumerate}
    \medskip
    \muted{The concept note later grows into your introduction. It is the first card on your board.}
  \end{columns}
\end{frame}

\begin{frame}{Project 1: foundation models for market risk}
  \framesubtitle{Task-specific heads on Chronos for volatility and tail-risk forecasting}
  \textbf{Idea.} Chronos models are time-series foundation models pretrained on large, mostly
  non-financial corpora. Keep the backbone and train lightweight \key{distributional or quantile
  heads} on its representations to forecast volatility, Value-at-Risk and Expected Shortfall.

  \smallskip
  \begin{columns}[T,onlytextwidth]
    \column{0.55\textwidth}
    \textbf{Questions}
    \begin{itemize}\small
      \item Does generic pretraining transfer to heavy-tailed, volatility-clustered returns?
      \item Zero-shot, new head on a frozen backbone, or fine-tuning: what is worth the compute?
      \item Do forecasts pass regulatory-style backtests, not only win on average loss?
    \end{itemize}
    \column{0.41\textwidth}
    \textbf{Benchmarks and evaluation}
    \begin{itemize}\small
      \item GARCH family, HAR-RV, quantile regression
      \item VaR coverage and independence tests
      \item joint VaR--ES scoring, model confidence sets
    \end{itemize}
  \end{columns}

  \muted{\small Splits naturally: volatility $\mid$ VaR and ES $\mid$ asset classes $\mid$ fine-tuning strategies.}
\end{frame}

\begin{frame}{Project 2: detrending and deseasonalisation}
  \framesubtitle{A review and benchmark: what to do about non-stationarity before modelling a time series}
  \textbf{Idea.} Every forecasting pipeline makes a choice about trend and seasonality, usually by
  habit. Build a \key{systematic review and a controlled benchmark} of these choices.

  \medskip
  \begin{columns}[T,onlytextwidth]
    \column{0.48\textwidth}
    \textbf{Methods to map}
    \begin{itemize}\small
      \item differencing and transformations
      \item filters: Hodrick--Prescott, Hamilton, band-pass
      \item decompositions: STL, MSTL, X-13, TRAMO-SEATS
      \item Fourier terms, wavelets, EMD, VMD
      \item learned: RevIN, neural decomposition blocks
    \end{itemize}
    \column{0.48\textwidth}
    \textbf{Questions}
    \begin{itemize}\small
      \item When does preprocessing help, and when does it destroy signal or leak the future?
      \item Do ML, deep and foundation models still need it?
      \item How do conclusions change across macro, financial and retail demand series?
    \end{itemize}
  \end{columns}

  \medskip
  \muted{\small Splits naturally: review and taxonomy $\mid$ benchmark on public datasets $\mid$ economic case study.}
\end{frame}

\begin{frame}{Project 3: foundation models for tabular data}
  \framesubtitle{A review of pretrained models for tabular prediction, with economic applications}
  \textbf{Idea.} Gradient-boosted trees have ruled tabular data for a decade. Pretrained
  \key{tabular foundation models} (in-context learners such as the TabPFN line, LLM-based and
  cross-table approaches) now challenge them. Review the field and test the claims where
  economists need them.

  \medskip
  \textbf{Questions}
  \begin{itemize}
    \item In which regimes do they win: small samples, many categories, distribution shift?
    \item Are the predicted probabilities \key{calibrated} well enough for pricing and credit decisions?
    \item What about interpretability, cost and regulatory constraints in finance and insurance?
  \end{itemize}

  \smallskip
  \textbf{Applications.} Credit scoring, insurance pricing, churn and customer value.

  \smallskip
  \muted{\small Splits naturally: literature review $\mid$ benchmark vs.\ boosting $\mid$ calibration and interpretability.}
\end{frame}

% =============================================================================
\section{Writing standards}
% =============================================================================

\begin{frame}{The logic of a thesis}
  \begin{center}
  \begin{tikzpicture}[
      node distance=7mm,
      box/.style={draw=semRule,fill=semTint,rounded corners=2pt,minimum height=9mm,
                  minimum width=26mm,font=\small},
      arr/.style={-{Stealth[length=2.2mm]},semAccent,line width=0.9pt}]
    \node[box] (p) {Problem};
    \node[box,right=of p] (h) {Hypotheses};
    \node[box,right=of h] (m) {Methods, data};
    \node[box,right=of m] (s) {Structure};
    \draw[arr] (p) -- (h); \draw[arr] (h) -- (m); \draw[arr] (m) -- (s);
  \end{tikzpicture}
  \end{center}

  \begin{itemize}
    \item Each step follows from the previous one. If a chapter does not help answer the
          question, it does not belong.
    \item \key{Introduction}: problem, why it matters, hypotheses, methods and data, structure.
    \item \key{Every chapter}: what it does and why; the content; what was established and how it
          leads to the next chapter.
    \item \key{Data chapter}: how the data was generated and collected, exact meaning of variables,
          descriptive statistics, and what this implies for modelling.
    \item \key{Conclusions}: your results, limitations, practical use, future work.
    \item \key{Abstract}: understandable to someone who will never read the thesis.
  \end{itemize}
\end{frame}

\begin{frame}{Style: a professional document, not literature}
  \begin{columns}[T,onlytextwidth]
    \column{0.48\textwidth}
    \textbf{Text}
    \begin{itemize}
      \item Short sentences, precise words. No metaphors, no ``some'', ``many'', ``appropriate''.
      \item Claims need evidence: a citation or your own result. ``It is well known'' is neither.
      \item Cite what is relevant to \emph{your} question, not everything the paper says.
      \item Draw only the conclusions your results support. State limitations plainly.
    \end{itemize}
    \column{0.48\textwidth}
    \textbf{Figures and tables}
    \begin{itemize}
      \item Every one is referenced by number and \key{interpreted in the text}.
      \item The caption says what the reader should see, and the source is always given.
      \item Readable at print size; consistent notation, units and model names throughout.
      \item No screenshots of code or output. Tables are typeset, plots are vector graphics.
    \end{itemize}
  \end{columns}
\end{frame}

% =============================================================================
\section{Next steps}
% =============================================================================

\begin{frame}{Before our next meeting}
  \begin{enumerate}
    \item Create a GitHub account (if you do not have one) and send me your username.
    \item Read this deck, the second deck (\emph{From Idea to Article}) and the \texttt{docs/} folder.
    \item Install the toolchain: git, a TeX distribution or Overleaf, \texttt{uv} or \texttt{renv}.
    \item Build the \key{thesis template} from the repository once, to check that your setup works.
    \item Send your topic proposal: title, two or three sentences, data confirmation.
          Say whether you want to work alone, in a pair, or join a larger project.
    \item Book your weekly status slot on Google Meet.
  \end{enumerate}

  \bigskip
  \begin{block}{Everything is in one place}
    \centering\repo
  \end{block}
\end{frame}

\end{document}
